\documentclass{article}

\usepackage{booktabs}
\usepackage{array}
\usepackage{xcolor}
\usepackage{listings}
\usepackage[obeyspaces,spaces]{url}
\urlstyle{tt}

\definecolor{newcolor}{rgb}{0.85,0.30,0.00}
\lstset{basicstyle=\ttfamily\footnotesize, frame=single, keepspaces=true,
        showstringspaces=false, aboveskip=6pt, belowskip=4pt,
        moredelim=[is][\color{newcolor}]{@@}{@@}}

\title{Technical Requirements:\\ \texttt{core} and \texttt{parsing}\\[1mm]
       \large UKRAINIX WATERDRONE -- analysis-app}
\author{}
\date{17 November 2026}

\begin{document}
\maketitle

\section{Scope and conventions}

This document specifies only the packages \path{core} and \path{parsing}
of the analysis app (both inside \path{src/analysis_app/}). Nothing else of
the planned app is needed for this step, and nothing else should be written
yet.

\begin{itemize}
\item \textbf{No data validation.} Assume a well-formed CSV: header present,
      all columns numeric and of equal length, no NaN values, no duplicate
      timestamps, no GPS glitches. Pure analysis only.
\item \textbf{Data type.} All data are 1-D numpy arrays: \path{int64} for
      timestamps, \path{float64} for everything else. Angles in degrees,
      distances in metres at every interface.
\item \textbf{Style.} Plain functions, no classes (except the
      \path{NamedTuple} below); inputs are never modified in place. Each
      package's \path{__init__.py} re-exports its public functions. Python
      3.10 or newer is assumed.
\item \textbf{Done when} every example below reproduces (floating-point noise
      ignored, $x,y$ within 0.1~m), converting to $x,y$ and back returns the
      input within $10^{-9}$ degrees, and all functions work for any number of
      rows and parameter columns.
\end{itemize}

\section{Libraries to use}

Please use established libraries instead of hand-written replacements:

\begin{itemize}
\item \path{numpy} -- all arrays and all numerical work;
\item \path{pandas} -- reading the CSV (\path{pandas.read_csv}), then
      \path{.to_numpy()} for every column;
\item \path{pyproj} -- optional: cross-check of the coordinate conversion
      (azimuthal equidistant projection centred on the reference origin);
\item standard library -- \path{pathlib} for paths, \path{typing} for the
      \path{NamedTuple};
\item \path{pytest} -- optional: turn the examples below into tests.
\end{itemize}

\section{Files to create}

Exactly the seven files shown in \textcolor{newcolor}{\textbf{orange}} are to
be created; marked \texttt{[new]} as well. The two folders already exist.

\begin{lstlisting}
src/analysis_app/
|-- core/                           exists, holds only .gitkeep
|   |-- @@__init__.py@@                 [new]
|   |-- @@types.py@@                    [new]
|   |-- @@config.py@@                   [new]
|   \-- @@functional.py@@               [new]
\-- parsing/                        exists, holds only .gitkeep
    |-- @@__init__.py@@                 [new]
    |-- @@coordinates.py@@              [new]
    \-- @@csv_parser.py@@               [new]
\end{lstlisting}

\begin{itemize}
\item The two \path{__init__.py} files only re-export the public names of
      their package: \path{core} -- \path{Sample}, the column constants,
      \path{config_dir}, \path{samples_dir}, \path{read_dict_file},
      \path{pipe}, \path{map_parameters}; \path{parsing} -- all seven
      functions of Section~\ref{sec:parsing}.
\item Nothing else is created: no tests, no \path{pyproject.toml}, nothing in
      other packages. Once the files exist, the placeholder \path{.gitkeep}
      in each of the two folders can be removed.
\end{itemize}

\section{What is assumed to exist}

Nothing outside these five files is assumed; their dependencies are:

\begin{center}
\footnotesize
\begin{tabular}{@{}>{\raggedright\arraybackslash}p{0.28\textwidth}>{\raggedright\arraybackslash}p{0.22\textwidth}>{\raggedright\arraybackslash}p{0.36\textwidth}@{}}
\toprule
\textbf{File} & \textbf{Imports from project} & \textbf{Provides}\\
\midrule
\path{core/types.py} & -- & \path{Sample}, column-name constants\\
\path{core/config.py} & -- & \path{config_dir}, \path{samples_dir}, \path{read_dict_file}\\
\path{core/functional.py} & -- & \path{pipe}, \path{map_parameters}\\
\path{parsing/coordinates.py} & -- & \path{choose_reference_origin}, \path{latlon_to_xy}, \path{xy_to_latlon}\\
\path{parsing/csv_parser.py} & \path{core.types}, \path{parsing.coordinates} & \path{read_columns}, \path{parameter_names}, \path{build_sample}, \path{load_sample}\\
\bottomrule
\end{tabular}
\end{center}

% =====================================================================
\section{Package \texttt{core}}

\subsection{\texttt{core/types.py}}

Constants and one data container; no other containers are needed yet.

\begin{lstlisting}
COL_TIME, COL_LAT, COL_LON, COL_DEPTH = (
    "TIMESTAMP", "LATITUDE", "LONGITUDE", "Z")
FIXED_COLUMNS = (COL_TIME, COL_LAT, COL_LON, COL_DEPTH)

class Sample(NamedTuple):
    t: np.ndarray        # (n,) int64   timestamps
    x: np.ndarray        # (n,) float64 metres, east of ref_origin
    y: np.ndarray        # (n,) float64 metres, north of ref_origin
    depth: np.ndarray | None   # (n,) float64, None if no Z column
    values: dict[str, np.ndarray]  # parameter -> (n,) float64
    ref_origin: tuple[float, float]  # (lat0, lon0) in degrees
\end{lstlisting}

\subsection{\texttt{core/config.py}}

\begin{itemize}
\item \path{config_dir() -> Path} -- the folder \path{analysis-app/config},
      found from the location of the file itself, not from the working
      directory.
\item \path{samples_dir() -> Path} -- the repository-level
      \path{samples/} folder, found the same way.
\item \path{read_dict_file(path) -> dict[str, str]} -- one entry per line in
      the form \path{KEY value}: key = first word, value = rest of the
      line, stripped. Empty lines and lines starting with \verb|#| are
      skipped. (This text format is a working assumption.)
\end{itemize}

\begin{lstlisting}
# file: "PARAM_A model_one" and "PARAM_B model_two" lines
>>> samples_dir()
PosixPath('/home/user/waterdrone-software/samples')
>>> read_dict_file(config_dir() / "physics_models.txt")
{'PARAM_A': 'model_one', 'PARAM_B': 'model_two'}
\end{lstlisting}

\subsection{\texttt{core/functional.py}}

\begin{itemize}
\item \path{pipe(value, *funcs)} -- applies the functions left to right:
      \path{funcs[-1](...funcs[0](value))}. With no functions it returns
      \path{value}.
\item \path{map_parameters(func, sample) -> dict} -- calls
      \path{func(sample, name)} once for every parameter name in
      \path{sample.values} (file order) and returns \verb|{name: result}|.
\end{itemize}

\begin{lstlisting}
>>> pipe(np.array([1., 4., 9.]), np.sqrt, np.sum)
6.0
>>> # sample: the one built in Section 6; result shown rounded
>>> map_parameters(lambda s, p: s.values[p].mean(), sample)
{'PARAM_A': 12.7, 'PARAM_B': 0.8067}
\end{lstlisting}

% =====================================================================
\section{Package \texttt{parsing}}
\label{sec:parsing}

\subsection{\texttt{parsing/coordinates.py}}

Converts GPS positions to local Cartesian coordinates: $x$ towards east, $y$
towards north, in metres, relative to a reference origin.

\begin{itemize}
\item \path{choose_reference_origin(lat, lon) -> (lat0, lon0)} -- centre of
      the bounding box of the points. Called once, for the first sample; later
      samples reuse the same origin so that all samples overlay.
\item \path{latlon_to_xy(lat, lon, ref_origin) -> (x, y)} -- arrays in,
      arrays of the same shape out; the reference origin maps to $(0,0)$.
      Use the method from the team's lecture notes \emph{T3 -- geographic
      coordinates translation}. A small-area approximation is sufficient:
      $x=R\cos\varphi_0\,(\lambda-\lambda_0)$,
      $y=R\,(\varphi-\varphi_0)$ with angles in radians and
      $R=6\,371\,000$~m. The examples below use this formula; any reasonable
      projection agrees with them to within 0.1~m.
\item \path{xy_to_latlon(x, y, ref_origin) -> (lat, lon)} -- inverse of
      \path{latlon_to_xy}. A round trip returns the input to within
      $10^{-9}$ degrees.
\end{itemize}

\noindent\begin{minipage}{\textwidth}
\begin{lstlisting}
>>> lat = np.array([50.0000, 50.0001, 50.0002])
>>> lon = np.array([30.0000, 30.0002, 30.0004])
>>> ref = choose_reference_origin(lat, lon)
>>> ref
(50.0001, 30.0002)
>>> x, y = latlon_to_xy(lat, lon, ref)       # shown rounded
>>> x
array([-14.295,   0.   ,  14.295])
>>> y
array([-11.119,   0.   ,  11.119])
>>> lat2, lon2 = xy_to_latlon(x, y, ref)
>>> np.allclose(lat2, lat), np.allclose(lon2, lon)
(True, True)
\end{lstlisting}
\end{minipage}

\subsection{\texttt{parsing/csv\_parser.py}}

\begin{itemize}
\item \path{read_columns(path) -> dict[str, ndarray]} -- reads the CSV with
      its header row; returns \verb|{column name: 1-D array}| in file order,
      \path{TIMESTAMP} as \path{int64}, all other columns as
      \path{float64}.
\item \path{parameter_names(columns) -> list[str]} -- all column names that
      are not in \path{FIXED_COLUMNS}, in file order.
\item \path{build_sample(columns, ref_origin=None) -> Sample} -- computes
      \path{x, y} from \path{LATITUDE} and \path{LONGITUDE} with
      \path{latlon_to_xy}; if \path{ref_origin} is \path{None} it first
      calls \path{choose_reference_origin}. \path{depth} is the \path{Z}
      column, or \path{None} if there is none. \path{values} holds one array
      per parameter name.
\item \path{load_sample(path, ref_origin=None) -> Sample} --
      \path{read_columns} followed by \path{build_sample}.
\end{itemize}

Example file \path{lake_run_03.csv}:

\begin{lstlisting}
TIMESTAMP,LATITUDE,LONGITUDE,Z,PARAM_A,PARAM_B
1730500000,50.00000,30.00000,0.5,12.4,0.83
1730500001,50.00010,30.00020,0.5,12.6,0.80
1730500002,50.00020,30.00040,0.5,13.1,0.79
\end{lstlisting}

\begin{lstlisting}
>>> cols = read_columns("lake_run_03.csv")
>>> list(cols)
['TIMESTAMP', 'LATITUDE', 'LONGITUDE', 'Z', 'PARAM_A', 'PARAM_B']
>>> parameter_names(cols)
['PARAM_A', 'PARAM_B']
>>> sample = load_sample("lake_run_03.csv")
>>> sample.ref_origin
(50.0001, 30.0002)
>>> sample.x                                 # shown rounded
array([-14.295,   0.   ,  14.295])
>>> sample.y
array([-11.119,   0.   ,  11.119])
>>> sample.values["PARAM_B"]
array([0.83, 0.8 , 0.79])
>>> # a second file in the same local frame as the first:
>>> other = load_sample("another.csv", ref_origin=sample.ref_origin)
\end{lstlisting}

\end{document}
